
It is readily seen, by reduction to the affine case and to an
algebraically closed base field, that if
\[
  u\colon G\longrightarrow G'
\]
is an epimorphism of smooth algebraic groups, then \(u(B)=B'\) is a
Borel subgroup of \(G'\) for every Borel subgroup \(B\) of \(G\).
We now consider the case in which this gives a bijective
correspondence between the Borel subgroups of \(G\) and those of
\(G'\).

\paragraph{\textbf{Proposition 4.9.}}
\SGARefTarget{sga3:xiv:4.9}
{\itshape Let \(G\) and \(G'\) be smooth \(S\)-preschemes in groups
of finite presentation with connected fibers, and let
\[
  u\colon G\longrightarrow G'
\]
be a faithfully flat, that is, surjective, homomorphism of
groups.\footnote{The editors ask whether a reference should be
added here.} Put \(N=\ker u\), and suppose that one of the following
two conditions holds:

\smallskip
\noindent
\textup{a)} \(N\) is central in \(G\).

\smallskip
\noindent
\textup{b)} \(S=\operatorname{Spec}(k)\) for a field \(k\), and, if
\(\bar k\) is an algebraic closure of \(k\), then
\[
  N_{\bar k}=\ker(u_{\bar k})
\]
is contained in the radical of \(G_{\bar k}\), that is, in its
largest smooth, connected, solvable, normal subgroup.

\smallskip
Then the map
\[
  B'\longmapsto u^{-1}(B')
\]
induces a bijection from the set of Borel subgroups of \(G'\) onto
the set of Borel subgroups of \(G\).}

Case \textup{b)} follows immediately from the correspondence between
algebraic subgroups of \(G'\) and algebraic subgroups of \(G\)
containing \(N\), together with the fact that, over an algebraically
closed field, the Borel subgroups of \(G\) contain its radical. The
latter fact follows at once from the argument preceding \SGARefLink{sga3:xiv:4.2}{4.2}.

We prove case \textup{a)}. By \SGARefLink{sga3:xii:7.12}{XII~7.12}, the map
\[
  H'\longmapsto H=u^{-1}(H')
\]
gives a bijection between the subpreschemes in groups \(H'\) of
\(G'\) that are smooth and of finite presentation over \(S\), have
connected fibers, and at every \(s\in S\) have the same reductive
rank and the same nilpotent rank as \(G'\), and the subpreschemes
in groups \(H\) of \(G\) having the analogous properties. Borel
subgroups of \(G'\) and of \(G\) have these properties.

It remains to show that, when \(H'\) and \(H\) correspond, \(H'\)
is a Borel subgroup of \(G'\) if and only if \(H\) is a Borel
subgroup of \(G\). By definition, this question reduces to the case
where \(S\) is the spectrum of an algebraically closed field. Since
\(N\) is central in \(G\), and hence in \(H\), the group \(H\) is
solvable if and only if \(H'\) is solvable. Finally, the
correspondence between algebraic subgroups of \(G'\) and algebraic
subgroups of \(G\) containing \(N\) shows immediately that \(H'\)
has the maximality property of \SGARefLink{sga3:xiv:4.1}{Definition~4.1} if and only if \(H\)
does. This completes the proof.

\paragraph{\textbf{Corollary 4.10.}}
\SGARefTarget{sga3:xiv:4.10}
{\itshape With the notation of \SGARefLink{sga3:xiv:4.9}{4.9},\footnote{The source prints
``\SGARefLink{sga3:xiv:4.9}{4.7}''; the intended reference is \SGARefLink{sga3:xiv:4.9}{4.9}.} if \(B'\) and \(B\) are
corresponding Borel subgroups of \(G'\) and \(G\), then
\[
  \mathfrak b=\operatorname{Lie}(u)^{-1}(\mathfrak b'),
\]
where \(\mathfrak g,\mathfrak g',\mathfrak b,\mathfrak b'\) are
the Lie algebras of \(G,G',B,B'\), respectively, and
\[
  \operatorname{Lie}(u)\colon\mathfrak g\longrightarrow\mathfrak g'
\]
is the homomorphism induced by \(u\).}

This follows immediately from the definitions and the identity
\[
  u^{-1}(B')=B.
\]

We can now prove the principal result of this section.
